\begin{table}[t]
\centering
\caption{Agreement Beyond Cohen's $\kappa$}
\label{tab:agreement-extended}
\footnotesize
\begin{tabular}{@{}L{3.8cm}rrrrr@{}}
\toprule
Measure & Agree (\%) & $\pm$1 (\%) & $\kappa$ / $\kappa_w$ & AC1 & $\alpha$ \\
\midrule
Hallucination & 94.7 & -- & 0.53 & 0.94 & 0.53 \\
Correctness & 88.0 & 96.0 & 0.63 & -- & 0.53 \\
Grade fit & 77.3 & 97.3 & 0.89 & -- & 0.69 \\
Clarity & 77.3 & 99.3 & 0.79 & -- & 0.81 \\
Restraint & 66.7 & 97.3 & 0.85 & -- & 0.75 \\
\midrule
MathDial: mistake identified & 96.0 & -- & 0.92 & 0.92 & 0.92 \\
\quad judge vs.\ member (lower) & 94.0 & -- & 0.88 & 0.88 & 0.88 \\
MathDial: answer revealed & 100.0 & -- & 1.00 & 1.00 & 1.00 \\
\quad judge vs.\ member (lower) & 96.0 & -- & 0.92 & 0.92 & 0.92 \\
MathDial: guidance & 24.0 & 86.0 & 0.47 & -- & 0.18 \\
\quad judge vs.\ member (lower) & 34.0 & 86.0 & 0.35 & -- & 0.30 \\
\bottomrule
\end{tabular}
\par\smallskip\parbox{\linewidth}{\scriptsize Two team members, 150 blind rows (MathDial: 50). $\kappa$ for yes/no labels, quadratic-weighted $\kappa_w$ for 1--5 scores. AC1: Gwet's, for yes/no labels. $\alpha$: Krippendorff's, nominal for yes/no and ordinal for scores. When nearly every label is the same (hallucination: 95\% agreement), $\kappa$ is low while AC1 stays high.}
\end{table}
